%% ============================================================
%%  Résumé (français)
%% ============================================================
\chapter*{Résumé}
\addcontentsline{toc}{chapter}{Résumé}
\thispagestyle{plain}

Le trading quantitatif sur la Bourse de Casablanca (MASI) souffre de deux problèmes
structurels: l'identification rapide des meilleures opportunités de marché dans un
environnement informationnel dense, et la validation rigoureuse des stratégies sans
overfitting. Ce projet de fin d'études, réalisé au sein de la Salle des Marchés de
\textbf{BMCE Capital}, répond à ces deux enjeux par la conception et le développement
d'une plateforme quantitative complète structurée en quatre stations.

\medskip

La plateforme implémente deux moteurs de résultat complémentaires. Le premier, le
\textit{Signal Engine}, évalue 20 familles d'indicateurs techniques sur 73 tickers MASI
via un pipeline de sept couches (Layers A→G) reposant sur l'évaluation par sous-périodes fixes
par fenêtres glissantes, un score de robustesse multi-composante et une réduction
de redondance par clustering de corrélation. Le second, le \textit{moteur de Backtest
WFO}, applique la méthode de Walk-Forward Analysis de Pardo (2008) avec la fonction
objectif PROM, produisant un verdict de robustesse via le Walk-Forward Efficiency (WFE).

\medskip

Ces deux moteurs alimentent un \textbf{tableau de bord unifié} permettant aux traders
de passer de la lecture du marché à la décision d'investissement — direction, conviction,
niveaux actionnables — en moins de 30 secondes par titre.

\medskip

Techniquement, la plateforme repose sur une stack full-stack (Next.js, FastAPI, PostgreSQL,
MinIO, RQ) avec un noyau quantitatif optimisé par Numba JIT atteignant un speedup de 38×
sur le kernel de backtesting.

\bigskip
\noindent\textbf{Mots-clés :}
backtesting quantitatif, walk-forward analysis, signal engine, indicateurs techniques,
évaluation hors-échantillon, MASI, trading algorithmique, overfitting, PROM, Numba.
